\begin{pdSolution}{sealed:ex:sealed-pairs}
Both parity classes have four members, so the equal-parity pairs number
$2\binom{4}{2}=12$; those must stay indistinguishable forever. The remaining
$\binom{8}{2}-12=16$ pairs have different parities and may differ at explicit
gate-release steps only. Table~\ref{sealed:tab:sealed-parity-pairs} shows how this
scales: doubling the secret count roughly quadruples both counts, since
both grow quadratically in $n$ while the equal-parity share of the total
stays fixed at one half.
\end{pdSolution}
\begin{pdSolution}{sealed:ex:sealed-m1-trace}
The script prints
\begin{quote}\footnotesize\texttt{[m1 leaky gate: releases raw s] caught: EQUAL-CLASS VIOLATION}\\
\texttt{witness: secrets (0,2), trace load\_data -> read\_secret ->}\\
\texttt{submit\_to\_gate -> gate\_release, Erin sees (('gate', 0),) vs (('gate', 2),)}
\end{quote}
The pair $(0,1)$ has different parities, so the honest gate is
\emph{allowed} to let Erin's view differ at a release step; a difference there
proves nothing. The pair $(0,2)$ has equal parity, so any difference is a
violation, and it appears at the fourth step, the release, when the leaky gate
emits the raw secret instead of its parity.
\end{pdSolution}
\begin{pdSolution}{sealed:ex:sealed-bypass}
The witness is \texttt{load\_data -> read\_secret -> BYPASS\_write\_E}, with
Erin seeing \texttt{(('leak', 0),)} against \texttt{(('leak', 1),)}: complete
mediation fails, since a release occurred with no gate decision. The two-run
check sees it because the bypass write lands in Erin-observable state at a step
that is not a gate release, which is a difference the property forbids for
every pair, whatever their parity.
\end{pdSolution}
\begin{pdSolution}{sealed:ex:sealed-launder}
Add a worker register $t$, a transition $\mathtt{compute}(f)$ that sets
$t=f(s)$ for $f$ drawn from a finite family, and let \texttt{submit} pass $t$
to the gate, which releases $g(t)$. The property must then be stated over
worlds that agree on $g(s)$ but may disagree on $g(f(s))$, and the checker must
quantify over the family of $f$; for any non-constant $f$ that maps an
equal-parity pair to a different-parity pair, the honest gate releases a
difference and the attack is exhibited. Delimited release rules the attack out
by requiring that no variable used in a declassification is updated before it,
which a compiler can check because the declassified expression is syntactic.
A language model's release candidate is arbitrary output, not an expression
the checker can identify, so the condition cannot be enforced and the residual
is priced as channel capacity instead.
\end{pdSolution}
\begin{pdSolution}{sealed:ex:sealed-compound}
The compound rule disables only \texttt{net\_egress}, a controllable event, in
the tainted state, and never disables the uncontrollable read, so its closed
loop is controllable and a supervisor realises it exactly. The second rule
must disable \texttt{in\_context\_read} itself, an event in $\Sigma_u$ that the
plant enables, so controllability fails on the first history that reads and the
rule is detect-only.
\end{pdSolution}
\begin{pdSolution}{sealed:ex:sealed-torn}
Under the torn write each release logs $\varepsilon$ but adds $\varepsilon-1$.
Client 0 releases 2: log 2, $\sigma=1$, $\sum_\Lambda=2$. Client 0 releases 2
again: the guard sees $\sigma+2=3\le4$ and admits it; log 2, $\sigma=2$,
$\sum_\Lambda=4$. Client 1 releases 1: the guard sees $\sigma+1=3\le4$ and
admits it; log 1, $\sigma=2$, $\sum_\Lambda=5>4$. The script prints exactly
this trace:
\begin{quote}\footnotesize\texttt{c0.release(2) [TORN: log 2, sigma +1] ->}\\
\texttt{c0.release(2) [TORN: log 2, sigma +1] -> c1.release(1) [TORN: log 1, sigma +0]}
\end{quote}
The balance never breaks the bound, so the guard keeps admitting, and the
auditable log overspends.
\end{pdSolution}
\begin{pdSolution}{sealed:ex:sealed-15-states}
A denied release still advances the denying client's program counter, so the
state records not only what was committed but which requests have already been
refused. Two interleavings that commit the same releases but differ in whether
a client's later request has already been refused are distinct states with the
same multiset. For example, after client 0 commits both of its releases
($\sigma=4$), the state before client 1's first request is refused and the
state after it is refused carry the same ledger $\{2,2\}$ and different
counters.
\end{pdSolution}
\begin{pdSolution}{sealed:ex:sealed-crossover}
At $k=8$, $\varepsilon=0.5$: basic $4.00$; advanced
$\sqrt{16\ln 10^{6}}\times0.5+8\times0.5\times(e^{0.5}-1)=7.43+2.59=10.03$,
so quote the sequential bound, by a wide margin [verified]. At
$\varepsilon=0.1$ the advanced bound is
$0.1\sqrt{2k\ln 10^{6}}+0.010517k$; at $k=34$ it is $3.42$ against basic
$3.40$, and at $k=35$ it is $3.48$ against $3.50$, so the crossover is $k=35$
[verified].
\end{pdSolution}
\begin{pdSolution}{sealed:ex:sealed-power}
$\beta^5=0.00032$, so detection is $0.99968$ [verified]; the script's analytic
line reads \texttt{k=5: analytic 0.99968 simulated 0.99967}. The approximation
gives $1-(1-0.01\times0.8)^{100}=1-0.992^{100}\approx0.552$.
The exact hypergeometric probability is approximately $0.554$; before rounding,
the approximation is lower by about $0.00143$\pdprov{a4\_canary\_sprt.py}{20260816}{internal}.
\end{pdSolution}
\begin{pdSolution}{sealed:ex:sealed-sprt}
% Keep the calculation and its essential statistical qualifications together.
\interlinepenalty=10000\relax
$\mathrm{KL}(p_1\Vert p_0)=0.01\ln10+0.99\ln(0.99/0.999)\approx0.01407$.
The numerator is $0.95\ln95+0.05\ln(0.05/0.99)\approx4.177$;
their quotient is $296.9$ [verified arithmetic]. Boundary values replace terminal
log-likelihoods; nominal targets replace attained errors. Overshoot, sampling
uncertainty and the cap can affect the reported $359$, a capped sample mean, not an
exact expectation. Non-detection
$0.0483$ and false alarm $0.0052$ are point estimates, not upper bounds.
Stopping counts include either decision; no wall-clock arrival model is specified.
\end{pdSolution}
\begin{pdSolution}{sealed:ex:sealed-canary-secrecy}
The clause: canary secrecy is a named security boundary, so disclosure of the
planting list is a breach with its own bond, and the list is regenerated per
job. The mechanism: plant per job from a fresh seed held by Derek's gateway
alone, so a stolen list has value for one job only and the compromise
probability applies per job rather than per engagement; power then returns
toward $1-\beta^k$ for every job whose list was not taken. What no redesign
restores is power against an adversary who holds \emph{this} job's list: with
the canaries stripped the leak carries none and Theorem~\ref{sealed:thm:sealed-power}
gives $1-\beta^0=0$, which is why detection is priced as a bounded bond and not
sold as a guarantee.
\end{pdSolution}
\begin{pdSolution}{sealed:ex:sealed-budget}
Before timing, $q\cdot b=200\times64=12{,}800$ bits. Three timing slots add up
to $\log_2 3\approx1.585$ bits per job, so about $317$ more bits, for roughly
$13{,}117$ bits in total; the timing term is why scheduled release with one
slot is the design default.
\end{pdSolution}
